\documentclass[12pt,a4paper]{article}

\usepackage[utf8]{inputenc}
\usepackage[T1]{fontenc}
\usepackage[brazil]{babel}
\usepackage{geometry}
\usepackage{booktabs}
\usepackage{tabularx}
\usepackage{longtable}
\usepackage{xcolor}
\usepackage{fancyhdr}
\usepackage{titlesec}
\usepackage{parskip}
\usepackage{microtype}
\usepackage{listings}
\usepackage{hyperref}
\usepackage{array}
\usepackage{lmodern}

\geometry{
  a4paper,
  top=2.5cm, bottom=2.5cm,
  left=3cm, right=2.5cm,
  headheight=15pt
}

% Cores
\definecolor{bhblue}{RGB}{0, 82, 147}
\definecolor{bhgray}{RGB}{90, 90, 90}
\definecolor{codebg}{RGB}{245, 245, 245}
\definecolor{codeframe}{RGB}{200, 200, 200}

% Hiperlinks
\hypersetup{
  colorlinks=true,
  linkcolor=bhblue,
  urlcolor=bhblue,
  citecolor=bhblue,
  pdftitle={Opensource BH — Documentação do Projeto},
  pdfauthor={paraenseembh},
  pdfsubject={Legislação Municipal de Belo Horizonte}
}

% Títulos
\titleformat{\section}
  {\Large\bfseries\color{bhblue}}
  {\thesection.}{0.6em}{}[\vspace{2pt}\color{bhblue}\rule{\linewidth}{0.8pt}]

\titleformat{\subsection}
  {\large\bfseries\color{bhgray}}
  {\thesubsection.}{0.5em}{}

\titleformat{\subsubsection}
  {\normalsize\bfseries}
  {\thesubsubsection.}{0.4em}{}

% Código
\lstset{
  backgroundcolor=\color{codebg},
  frame=single,
  rulecolor=\color{codeframe},
  basicstyle=\ttfamily\small,
  breaklines=true,
  breakatwhitespace=true,
  columns=fullflexible,
  keepspaces=true,
  showstringspaces=false,
  upquote=true,
  xleftmargin=8pt,
  xrightmargin=8pt,
  aboveskip=8pt,
  belowskip=8pt,
  extendedchars=false,
}

% Cabeçalho e rodapé
\pagestyle{fancy}
\fancyhf{}
\fancyhead[L]{\small\color{bhgray}Opensource BH}
\fancyhead[R]{\small\color{bhgray}Documentação do Projeto}
\fancyfoot[C]{\small\thepage}
\renewcommand{\headrulewidth}{0.4pt}

% Sem recuo de parágrafo, espaço entre parágrafos
\setlength{\parindent}{0pt}
\setlength{\parskip}{6pt}

% ----------------------------------------------------------------
\begin{document}

% Página de título
\begin{titlepage}
  \centering
  \vspace*{3cm}
  {\Huge\bfseries\color{bhblue} Opensource BH\par}
  \vspace{0.5cm}
  {\Large\color{bhgray} Documentação do Projeto\par}
  \vspace{2cm}
  \rule{0.6\linewidth}{1pt}\par
  \vspace{1.5cm}
  {\large
    Ferramentas abertas para consulta e exploração\\
    da legislação municipal de \textbf{Belo Horizonte}
  \par}
  \vspace{3cm}
  {\normalsize
    Repositório: \href{https://github.com/paraenseembh/opensource_bh}%
    {github.com/paraenseembh/opensource\_bh}\par
  }
  \vfill
  {\small\color{bhgray} Maio de 2026}
\end{titlepage}

\tableofcontents
\newpage

% ================================================================
\section{O que é este projeto?}

Este repositório reúne ferramentas abertas para consultar e explorar a legislação municipal de \textbf{Belo Horizonte}. O foco é democratizar o acesso a decretos, portarias e leis municipais, tornando esses documentos pesquisáveis e compreensíveis por qualquer pessoa --- não apenas por advogados ou servidores públicos.

O projeto nasce de uma ideia simples: os atos normativos da prefeitura estão publicados na internet, mas dispersos, difíceis de buscar e escritos em linguagem técnica. Com um chatbot e um motor de busca, qualquer cidadão pode perguntar em linguagem natural --- \textit{"quais decretos tratam de transporte público em 2026?"} --- e obter uma resposta compreensível.

% ================================================================
\section{Módulos do repositório}

\begin{tabularx}{\linewidth}{lX}
  \toprule
  \textbf{Pasta} & \textbf{O que faz} \\
  \midrule
  \texttt{bh\_news\_chatbot/}    & Chatbot em Python que responde perguntas sobre a legislação municipal usando IA \\
  \texttt{bh\_news\_chatbot\_r/} & A mesma ferramenta, reescrita em R para quem prefere essa linguagem \\
  \texttt{corpus\_dom/}          & Pipeline de mineração: baixa e estrutura documentos do Diário Oficial do Município (DOM-BH) \\
  \texttt{busca\_dom/}           & Motor de busca com interface web para pesquisar no corpus do DOM-BH \\
  \texttt{legislacao\_2026\_04.csv} & Base de dados com 5.197 registros de decretos, portarias e leis (abril de 2026) \\
  \bottomrule
\end{tabularx}

% ================================================================
\section{Principais marcos do desenvolvimento}

A tabela abaixo registra os commits mais relevantes, em ordem cronológica --- do mais antigo ao mais recente.

\begin{longtable}{lp{10.5cm}}
  \toprule
  \textbf{Commit} & \textbf{O que foi feito} \\
  \midrule
  \endhead
  \texttt{1bb853f} & Criação do chatbot inicial, lendo links de legislação a partir de uma planilha Google Sheets \\[4pt]
  \texttt{0917368} & Adição de categorização automática de decretos por tema (urbanismo, saúde, educação etc.) via IA \\[4pt]
  \texttt{f90c78f} & Suporte ao Google Gemini como provedor de IA, além do Claude (Anthropic) \\[4pt]
  \texttt{6d82c97} & Versão do chatbot reescrita em linguagem R \\[4pt]
  \texttt{1b661a1} & Correção do acesso ao Google Sheets (API descontinuada substituída) \\[4pt]
  \texttt{601478b} & Adição do pipeline \texttt{corpus\_dom}: extração e estruturação dos PDFs e páginas do DOM-BH \\[4pt]
  \texttt{48811ea} & Suporte a PDF e DOCX no scraper; exportação para JSONL e CSV \\[4pt]
  \texttt{22c45ab} & Motor de busca \texttt{busca\_dom} com TF-IDF e interface Streamlit \\[4pt]
  \texttt{c0b7bca} & Upload do CSV de legislação com mais de 5.000 registros \\[4pt]
  \texttt{041ec3f} & Adaptação do projeto para foco nos decretos municipais de BH \\[4pt]
  \texttt{8107507} & Adição da Maritaca AI (modelo Sabiá, IA brasileira) como terceiro provedor \\[4pt]
  \texttt{a5a671d} & Script de verificação das APIs dos três provedores \\[4pt]
  \texttt{5647369} & Testes dedicados para a Maritaca AI com cobertura de erros \\[4pt]
  \texttt{b120473} & Atualização da documentação principal do README \\
  \bottomrule
\end{longtable}

% ================================================================
\section{Dependências}

\subsection{Python --- Chatbot (\texttt{bh\_news\_chatbot/})}

\begin{lstlisting}
anthropic>=0.40.0       # API do Claude (Anthropic)
google-genai>=1.0.0     # API do Google Gemini
openai>=1.0.0           # Compatibilidade OpenAI (usada pela Maritaca)
beautifulsoup4>=4.12.0  # Raspagem de paginas web
lxml>=5.0.0             # Parser HTML/XML
requests>=2.32.0        # Requisicoes HTTP
\end{lstlisting}

\subsection{Python --- Corpus DOM (\texttt{corpus\_dom/})}

\begin{lstlisting}
pandas>=2.0.0                   # Manipulacao de dados tabulares
requests>=2.31.0                # Requisicoes HTTP
beautifulsoup4>=4.12.0          # Raspagem de HTML
lxml>=4.9.0                     # Parser HTML/XML
mysql-connector-python>=8.3.0   # Conexao com banco de dados MySQL
python-dotenv>=1.0.0            # Leitura de variaveis de ambiente
pdfplumber>=0.10.0              # Extracao de texto de PDFs
python-docx>=1.0.0              # Leitura de arquivos DOCX
\end{lstlisting}

\subsection{Python --- Busca DOM (\texttt{busca\_dom/})}

\begin{lstlisting}
pandas>=2.0.0           # Manipulacao de dados
scikit-learn>=1.4.0     # Algoritmo de busca TF-IDF
streamlit>=1.35.0       # Interface web interativa
python-dotenv>=1.0.0    # Leitura de variaveis de ambiente
\end{lstlisting}

\subsection{R --- Chatbot (\texttt{bh\_news\_chatbot\_r/})}

\begin{lstlisting}
httr2       # Requisicoes HTTP
rvest       # Raspagem de paginas web
jsonlite    # Leitura e escrita de JSON
readr       # Leitura de CSV
rlang       # Utilitarios de programacao em R
optparse    # Argumentos de linha de comando
\end{lstlisting}

% ================================================================
\section{Para quem não tem experiência com programação}

Esta seção é uma introdução gentil para quem quer rodar o projeto, mas nunca usou Git ou Python antes. Não se preocupe --- você não precisa entender tudo de uma vez. Cada passo é simples quando feito um de cada vez.

\subsection{O que é código aberto (opensource)?}

Um projeto \textbf{opensource} é aquele cujo código-fonte está disponível publicamente para qualquer pessoa ler, usar, modificar e distribuir. Pense como uma receita de bolo publicada na internet: qualquer um pode copiar, adaptar e melhorar.

Este projeto segue essa filosofia. Ele está disponível gratuitamente, aceita contribuições de qualquer pessoa e foi construído colaborativamente, usando ferramentas que também são abertas.

Plataformas como o \textbf{GitHub} hospedam projetos opensource e permitem que pessoas do mundo inteiro colaborem. O endereço deste projeto é \href{https://github.com/paraenseembh/opensource_bh}{\texttt{github.com/paraenseembh/opensource\_bh}}.

\subsection{O que é Git e por que usamos?}

\textbf{Git} é um sistema que registra todas as mudanças feitas no código ao longo do tempo. É como um histórico de versões, mas muito mais poderoso: você pode ver quem mudou o quê, quando, e por quê. Se algo der errado, dá para voltar a uma versão anterior.

Quando você ``clona'' um repositório, está baixando uma cópia completa do projeto, incluindo todo esse histórico.

\subsubsection{Instalando o Git}

\begin{itemize}
  \item \textbf{Linux (Ubuntu/Debian):} \lstinline|sudo apt install git|
  \item \textbf{Mac:} \lstinline|brew install git| (ou baixar em \href{https://git-scm.com}{git-scm.com})
  \item \textbf{Windows:} baixar o instalador em \href{https://git-scm.com/downloads}{git-scm.com/downloads}
\end{itemize}

\subsubsection{Clonando o repositório}

Abra o terminal e digite:

\begin{lstlisting}
git clone https://github.com/paraenseembh/opensource_bh.git
cd opensource_bh
\end{lstlisting}

Para baixar atualizações futuras sem clonar de novo:

\begin{lstlisting}
git pull
\end{lstlisting}

\subsubsection{Comandos Git úteis}

\begin{tabularx}{\linewidth}{lX}
  \toprule
  \textbf{Comando} & \textbf{O que faz} \\
  \midrule
  \texttt{git status}        & Mostra o que mudou desde o último registro \\
  \texttt{git log --oneline} & Exibe o histórico de commits de forma compacta \\
  \texttt{git pull}          & Baixa as últimas atualizações do projeto \\
  \texttt{git checkout X}    & Muda para a branch (versão) chamada X \\
  \bottomrule
\end{tabularx}

\subsection{O que é Python e como instalar?}

\textbf{Python} é uma linguagem de programação muito usada em ciência de dados, automação e inteligência artificial. O chatbot e os scripts deste projeto são escritos em Python.

Baixe a versão 3.10 ou mais recente em \href{https://python.org/downloads}{\texttt{python.org/downloads}}. No Windows, durante a instalação, marque a opção \textbf{``Add Python to PATH''}.

Verifique se funcionou:

\begin{lstlisting}
python --version
\end{lstlisting}

\subsubsection{Instalando as dependências}

As bibliotecas usadas pelo projeto estão listadas em \texttt{requirements.txt}. Para instalar tudo de uma vez:

\begin{lstlisting}
pip install -r bh_news_chatbot/requirements.txt
\end{lstlisting}

\subsubsection{Ambientes virtuais (boa prática)}

Um \textbf{ambiente virtual} isola as bibliotecas deste projeto sem interferir no resto do sistema:

\begin{lstlisting}
# Cria o ambiente (so precisa fazer uma vez)
python -m venv .venv

# Ativa - Linux/Mac
source .venv/bin/activate

# Ativa - Windows
.venv\Scripts\activate

# Instala as dependencias dentro do ambiente
pip install -r bh_news_chatbot/requirements.txt
\end{lstlisting}

Para desativar quando terminar:

\begin{lstlisting}
deactivate
\end{lstlisting}

\subsection{Rodando o chatbot}

Com as dependências instaladas e uma chave de API configurada:

\begin{lstlisting}
# Usando a base de dados local
python bh_news_chatbot/main.py --use-local-csv
\end{lstlisting}

Crie sua chave gratuita em um dos provedores suportados:

\begin{itemize}
  \item \textbf{Anthropic (Claude):} \href{https://console.anthropic.com}{console.anthropic.com}
  \item \textbf{Google Gemini:} \href{https://aistudio.google.com}{aistudio.google.com}
  \item \textbf{Maritaca AI (Sabiá --- IA brasileira):} \href{https://plataforma.maritaca.ai}{plataforma.maritaca.ai}
\end{itemize}

Configure a chave no terminal:

\begin{lstlisting}
# Linux/Mac
export ANTHROPIC_API_KEY=sua-chave-aqui

# Windows (CMD)
set ANTHROPIC_API_KEY=sua-chave-aqui
\end{lstlisting}

Ou passe diretamente na linha de comando:

\begin{lstlisting}
python bh_news_chatbot/main.py --use-local-csv --api-key sua-chave-aqui
\end{lstlisting}

Dentro do chat, pergunte em linguagem natural:

\begin{lstlisting}
> Quais decretos tratam de transporte publico em 2024?
> O que diz o decreto 18.012?
> Existem atos sobre habitacao de interesse social?
\end{lstlisting}

Use \texttt{/ajuda} para ver todos os comandos e \texttt{/sair} para encerrar.

\subsection{Como contribuir com o projeto}

Opensource é sobre colaboração. Se você encontrou um bug, quer sugerir uma melhoria ou adicionar uma funcionalidade:

\begin{enumerate}
  \item \textbf{Abra uma issue} no GitHub descrevendo o problema ou a ideia
  \item \textbf{Fork} o repositório (cria uma cópia sua no GitHub)
  \item Faça as mudanças na sua cópia
  \item \textbf{Abra um Pull Request} --- uma proposta de mudança que os mantenedores podem revisar e aceitar
\end{enumerate}

Não precisa ser programador para contribuir: melhorar a documentação, reportar bugs ou sugerir novas fontes de dados já são contribuições valiosas.

\vfill
\begin{center}
  \small\color{bhgray}
  \rule{0.4\linewidth}{0.4pt}\\[4pt]
  Opensource BH \textbullet{} Maio de 2026 \textbullet{}
  \href{https://github.com/paraenseembh/opensource_bh}{github.com/paraenseembh/opensource\_bh}
\end{center}

\end{document}
